% ============================================================
% Capítulo 07: Predicción de rendimientos y coeficientes de rendimiento
% ============================================================

% Consejo: Este capítulo se enfoca en cómo usar la estequiometría para predecir cuánto producto o biomasa se obtendrá de un sustrato.

\section{Predicción de rendimientos}

Los rendimientos son relaciones que nos indican la eficiencia con la que un sustrato se convierte en biomasa o producto. Se pueden predecir a partir de la ecuación estequiométrica balanceada del crecimiento microbiano.

\subsection{Rendimiento de biomasa respecto a sustrato ($Y_{X/S}$)}
Es la cantidad de biomasa producida por gramo de sustrato consumido.

\begin{equation}
    Y_{X/S} = \frac{\text{Biomasa producida (g)}}{\text{Sustrato consumido (g)}} = \frac{\Delta X}{\Delta S}
\end{equation}

\subsection{Rendimiento de producto respecto a sustrato ($Y_{P/S}$)}
Es la cantidad de producto producido por gramo de sustrato consumido.

\begin{equation}
    Y_{P/S} = \frac{\text{Producto producido (g)}}{\text{Sustrato consumido (g)}} = \frac{\Delta P}{\Delta S}
\end{equation}

% Consejo: En la práctica, estos rendimientos se obtienen experimentalmente. Sin embargo, el balance estequiométrico permite calcular rendimientos teóricos máximos.

\section{Coeficientes de rendimiento máximo teórico}

El \textbf{rendimiento máximo teórico} se calcula a partir de la estequiometría de la reacción, asumiendo que no hay formación de subproductos y que el sustrato se convierte completamente en biomasa o producto.

\subsection*{Ejemplo: Cálculo de rendimiento teórico de etanol a partir de glucosa}

La ecuación de la fermentación alcohólica es:

\[
C_6H_{12}O_6 \rightarrow 2 C_2H_5OH + 2 CO_2
\]

*   Masa molar de glucosa = 180 g/mol.
*   Masa molar de etanol = 46 g/mol.

Por cada 1 mol de glucosa (180 g), se producen 2 moles de etanol (2 * 46 = 92 g).

El rendimiento máximo teórico de etanol respecto a glucosa es:

\[
Y_{P/S} = \frac{92 \text{ g etanol}}{180 \text{ g glucosa}} = 0.511 \text{ g etanol / g glucosa}
\]

\section{Coeficientes de rendimiento real}

El \textbf{rendimiento real} es el que se obtiene experimentalmente en un bioproceso real.

\[
Y_{P/S}^{\text{real}} = \frac{P_{\text{obtenido}}}{S_{\text{consumido}}}
\]

Siempre es menor que el rendimiento máximo teórico debido a:
\begin{itemize}
    \item El uso de parte del sustrato para el mantenimiento celular.
    \item La producción de subproductos no deseados.
    \item La ineficiencia de los procesos metabólicos.
\end{itemize}